%% kapa: exam=1 round=23 year=2012 period=2 form=A subject=accounting subject_ko=회계학 seq=40 number=40
%% source: Q-Net 감정평가사 기출문제 (cst00309) artlSeq=5052580 file=raw/1008/5052580/2012년 제23회 감정평가사 제1차시험 기출문제(2교시 A형).pdf
%% conversion: manual transcription from page image; review=proofread-gpt-terra; flags=none
\begin{question}{40}
표준원가시스템을 사용하고 있는 (주)한국의 직접재료원가의 제품단위당 표준사용량은 10kg이고, 표준가격은 kg당 ￦6이다. (주)한국은 6월에 직접재료 40,000kg을 ￦225,000에 구입하여 36,000kg을 사용하였다. (주)한국의 6월 중 제품생산량은 3,000단위이다. 직접재료 가격차이를 구입시점에 분리하는 경우, 6월의 직접재료원가에 대한 가격차이와 능률차이(수량차이)는?
\begin{choices}
  \item 가격차이 ￦6,000불리, 능률차이 ￦32,000유리
  \item 가격차이 ￦9,000불리, 능률차이 ￦36,000불리
  \item 가격차이 ￦15,000유리, 능률차이 ￦36,000불리
  \item 가격차이 ￦15,000유리, 능률차이 ￦32,000불리
  \item 가격차이 ￦9,000불리, 능률차이 ￦36,000유리
\end{choices}
\end{question}
\answer{3}
